\markboth{STATE OF THE ART}{STATE OF THE ART}
\section{State of the Art}\label{sec:sota}

This section reviews the four research pillars Training Copilot builds upon: agentic AI and multi-agent systems, tool-integration protocols, sports analytics and wearable technology, and retrieval-augmented generation.

\subsection{Agentic AI and Multi-Agent Systems}\label{sec:sota:agents}

The rapid scaling of large language models has shifted the field from static text generation toward autonomous \textit{agentic} systems that plan, reason, and act in external environments. The Transformer architecture~\citep{vaswani_attention_2017}, chain-of-thought prompting~\citep{wei_chain-of-thought_2022}, and the scaling of model capacity to hundreds of billions of parameters~\citep{openai_gpt4_2023} laid the groundwork by showing that LLMs perform multi-step reasoning once prompted to generate intermediate steps. We organise the literature along three themes: reasoning architectures, agent-internal structure, and multi-agent coordination.

\subsubsection{Reasoning Architectures}

Two complementary patterns ground LLM reasoning in external action. \cite{yao_react_2023} introduced \textit{ReAct}, interleaving reasoning traces (``thought'' steps) with task-specific actions in a single loop: the model generates a thought, executes an action (e.g.\ \texttt{search[entity]} against a Wikipedia API), observes the result, and repeats until done. Evaluated on PaLM-540B, ReAct reached 71\,\% success on ALFWorld (vs.\ 45\,\% action-only) and 40\,\% on WebShop (vs.\ 30.1\,\%), and produced zero hallucinated reasoning chains on HotpotQA/FEVER versus 56\,\% for chain-of-thought---because every claim is grounded in a retrieved observation rather than generated from parameters. Training Copilot's specialists use this pattern directly: each is a LangGraph ReAct agent whose ``actions'' are MCP tool calls. \cite{shinn_reflexion_2023} extended the loop across episodes with \textit{Reflexion}, where agents store linguistic self-reflections after failures in an episodic buffer to inform later trials without weight updates---complementary to ReAct's within-episode loop. We use Reflexion in one subsystem: the LLM-driven chart generator retries once with the error message as feedback (Section~\ref{sec:architecture:supporting}). Both patterns, however, assume a single agent over a fixed tool set; neither addresses multi-agent coordination or dynamic tool discovery across heterogeneous sources.

\subsubsection{Agent-Internal Architecture}

\cite{park_generative_2023} showed that LLM-powered agents exhibit believable behaviour when equipped with a \textit{memory stream} recording experiences, a \textit{reflection} mechanism synthesising abstractions, and a \textit{planning} module translating reflections into action sequences; \cite{wang_survey_2024} generalised this into a four-module taxonomy---profiling, memory, planning, action---subsuming prior architectures. Both works find memory and planning mutually reinforcing (richer memory enables better planning, which yields more informative experience), but both describe single-agent architectures and do not address how agents with different profiles and tool sets should coordinate.

\subsubsection{Multi-Agent Coordination}

Three works illustrate the multi-agent design space. \cite{shen_hugginggpt_2023} proposed HuggingGPT, where a central controller decomposes requests, selects specialist models from a hub, executes them, and summarises results---demonstrating orchestrator--specialist decomposition, though applied to model selection rather than tool-based data retrieval. \cite{wu_autogen_2023} introduced AutoGen, showing empirically that role-specialised agents with tool access outperform monolithic designs on complex tasks, and that flexible conversation patterns (sequential, parallel, nested) enable different coordination strategies. \cite{guo_large_2024} surveyed the field and identified two recurring open challenges---communication protocol design and task decomposition---while \cite{acharya_agentic_2025} distinguished agentic AI from traditional automation by its capacity for autonomous goal pursuit through perception, reasoning, and action. These works converge on a shared finding: multi-agent architectures with domain-scoped specialists outperform monolithic agents on heterogeneous tasks, yet none demonstrate \textit{standardised} tool-discovery or inter-agent protocols, relying instead on ad-hoc registries and custom communication layers.

\subsection{Tool Integration Protocols}\label{sec:sota:protocols}

The practical utility of LLM agents depends on interacting with external tools and data. Research here progressed through three stages: teaching models to use tools at all, standardising how tools are discovered and invoked, and standardising how tool-using agents talk to each other.

\subsubsection{Tool-Augmented Language Models}

\cite{schick_toolformer_2023} showed that a 6.7B-parameter model (GPT-J) can teach itself to use tools through a self-supervised pipeline: candidate API calls are sampled within training text, executed, and kept only if they reduce the cross-entropy loss over subsequent tokens. Trained on five tools (calculator, QA, Wikipedia search, translation, calendar), the model selected the correct tool in 97.9\,\% of cases on mathematics benchmarks, outperforming GPT-3 (175B) despite being 25$\times$ smaller---though limited to one tool call per example with no query reformulation. \cite{patil_gorilla_2023} addressed tool selection at scale, fine-tuning LLaMA-7B on 16{,}450 instruction--API pairs across 1{,}645 APIs from three model hubs; on their APIBench benchmark, Gorilla beat GPT-4's zero-shot accuracy on TorchHub by 20.43 points and cut hallucinated API calls from 36.55\,\% (GPT-4) to 6.98\,\%, reaching 0\,\% on TorchHub with a ground-truth retriever. \cite{qu_tool_2024} synthesised these findings into a four-stage workflow---planning, selection, calling, response generation---and showed across benchmarks scaling to 16{,}464 tools (ToolBench) that placing all tool descriptions in context becomes impractical beyond a few hundred tools, making retrieval-based pre-filtering mandatory. This motivates Training Copilot's design directly: scoping each specialist to at most two MCP servers (1--26 tools per specialist) keeps tool selection well below that threshold.

\subsubsection{From Ad-hoc Integration to Standardised Protocols}

The proliferation of tool-augmented agents created an $M \times N$ fragmentation problem---$M$ applications each integrating $N$ data sources separately. Two complementary protocols address this. The \textbf{Model Context Protocol (MCP)}~\citep{anthropic_mcp_2024, anthropic_mcp_announcement_2024} defines a universal client--server protocol over JSON-RPC\,2.0~\citep{jsonrpc_spec_2013} for tool discovery and invocation, exposing Prompts, Resources, and Tools while separating authentication from tool declaration so a host can use arbitrary servers without provider-specific code. \cite{jayanti_enhancing_2026} proposed Context-Aware MCP (CA-MCP), adding a \textit{Shared Context Store} that servers read and write to after the central LLM seeds it with goals and constraints---reducing execution time by 67.8\,\% and LLM orchestration calls by 60\,\% on the TravelPlanner benchmark while improving task completeness. Training Copilot does not adopt this shared-memory pattern (A2A handles inter-agent coordination instead), but the result validates MCP's extensibility. The \textbf{Agent-to-Agent protocol (A2A)}~\citep{google_a2a_2025}, donated to the Linux Foundation after Google's April 2025 announcement, standardises inter-agent communication over JSON-RPC\,2.0 around \textit{opaque internals}: agents advertise capabilities via Agent Cards but expose only results, not reasoning, enabling composition across frameworks and providers. \cite{yang_survey_2025} distinguish the two as complementary: MCP is \textit{context-oriented} (agents to data), A2A is \textit{coordination-oriented} (agents to agents). No published system, however, has combined both protocols in a domain-specific multi-agent application.

\subsection{Sports Analytics and Wearable Technology}\label{sec:sota:sports}

Sports science converges on one finding central to this work: effective training management requires integrating multiple physiological and contextual signals, yet no consumer tool performs this integration with reasoning. We review three themes: training load and periodisation, recovery monitoring, and wearable reliability.

\subsubsection{Training Load: Monitoring and Injury Prevention}

\cite{bourdon_monitoring_2017}, in an international consensus statement (11 authors, five countries), formalise \textit{external load}---objective work performed, measured by GPS and accelerometers---against \textit{internal load}---the psychophysiological response (heart rate, HRV, session RPE, lactate)---and recommend monitoring both jointly, since identical external load produces different internal responses depending on fatigue or illness; this ``uncoupling'' is itself diagnostic. They establish the acute-to-chronic workload ratio (ACWR), with 0.8--1.3 as low-risk; \cite{halson_monitoring_2014} add that no single biomarker is sufficiently validated alone. \cite{gabbett_training_2016} formalise the ``training--injury prevention paradox'' from cricket, Australian football, and rugby league data: well-conditioned athletes have \textit{fewer} injuries than undertrained ones, since conditioning itself protects. At ACWR $\leq$0.99, seven-day injury likelihood was $\sim$4\,\% in elite fast bowlers; at $\geq$1.5, risk rose 2--4$\times$, and week-to-week load increases of 15\,\%\,+ tracked with 21--49\,\% injury rates versus under 10\,\% when load stayed roughly constant week to week. \cite{grgic_periodization_2017} extend this to periodisation, finding in a meta-analysis that linear and undulating schemes produce equivalent hypertrophy---the specific periodisation pattern does not significantly affect outcomes. Together these findings establish that load management requires trend analysis over time, not point-in-time metrics; Training Copilot's load specialist computes CTL, ATL, and TSB from Strava data---derived from the fitness--fatigue impulse-response model of \cite{morton_modeling_1990}---to operationalise this. The ACWR itself, however, has faced methodological criticism: \cite{impellizzeri_training_2020} identify mathematical coupling (acute load appears in both numerator and denominator) and show that simulations can produce spurious ACWR--injury associations from unrelated data, so our system presents ACWR as one contextual signal rather than a standalone predictor.

\subsubsection{Recovery: From Single Metrics to Multi-Signal Assessment}

\cite{kellmann_recovery_2018} argue in a consensus statement that systematic recovery monitoring---not just load tracking---should be integral to training, since the stress--recovery balance determines long-term sustainability. HRV has emerged as a central biomarker, but its interpretation is not straightforward: \cite{plews_training_2013}, using data from Olympic and World Champion rowers, show that in elite athletes both increases \textit{and} decreases in HRV can signal negative adaptation, because vagal HRV saturates at very low resting heart rates (the HRV--R-R interval relationship is quadratic, not linear). They recommend weekly rolling averages of Ln\,rMSSD over single-day values (which showed no meaningful performance correlation, $r=-0.17$) interpreted against each athlete's own baseline---among four gold-medal rowers, the Ln\,rMSSD--R-R correlation ranged from $r=0.91$ to $r=-0.03$, profiles a fixed-threshold system would misread. \cite{lundstrom_hrv_2023} review HRV-guided training in consumer settings, noting the gap between clinical- and consumer-grade measurement. The strongest evidence for multi-signal approaches comes from \cite{rothschild_predicting_2024}, who tracked 43 endurance athletes (67\,\% triathletes) over 12 weeks (3{,}572 person-days) and tested nine ML algorithms predicting daily perceived recovery status: the best group model (LASSO) reached $R^2=0.31$ (RMSE 11.8 on 0--100), individual top-5-feature models reached $R^2=0.35\pm0.20$ with RMSE varying fivefold (5.5--23.6) across participants, and the most predictive features were soreness, sleep index, and prior-day status---not HRV alone. \cite{wackerhage_personalized_2021} frame this as inherent: a training plan is a complex mix of interacting, time-varying interventions, making purely evidence-based decisions impractical---motivating tool-assisted approaches like Training Copilot's recovery specialist, which queries HRV, sleep, Body Battery, and stress simultaneously to approximate this multi-signal assessment.

\subsubsection{Wearable Reliability and Data Ecosystem Biases}

Wearable data reliability is a prerequisite for any analytics system built on it. \cite{fuller_reliability_2020} reviewed 158 publications across 45 device models from nine manufacturers: of 979 step-count comparisons, 45\,\% of those with calculable error fell within $\pm$3\,\% of criterion measures (slight underestimation, median $-2\,\%$), and 57\,\% of heart-rate comparisons were within $\pm$3\,\% (Apple Watch 71\,\%, Fitbit 51\,\%, Garmin 49\,\%); precision degraded during high-intensity exercise and in free-living settings, energy expenditure was inaccurate across all devices, and while interdevice reliability for steps was very strong, intradevice reliability was variable (mean $r=0.58$). \cite{seckin_review_2023} note that translating raw sensor data into actionable coaching remains open, and \cite{vanhooren_wearables_2020} identify challenges in feedback timing and injury-risk communication for real-time wearable feedback in running. \cite{venter_bias_2023} further quantify selection bias in Strava's crowdsourced data---overrepresenting recreationally active, male, and middle-aged (35--54) users---a bias professional sports avoid~\citep{bourdon_monitoring_2017} but consumer platforms must acknowledge. Training Copilot mitigates this partially through multi-source corroboration (cross-referencing HRV, sleep, and Body Battery) but cannot fully compensate for sensor inaccuracy.

Across all three themes, sports science establishes that training decisions require integrating load trends, recovery signals, and contextual factors, yet the tools that \textit{collect} this data do not \textit{reason across} it---the gap is in synthesis, not availability.

\subsection{Retrieval-Augmented Generation}\label{sec:sota:rag}

Retrieval-Augmented Generation combines the factual precision of information retrieval with generative fluency. \cite{lewis_retrieval-augmented_2020} introduced the foundational architecture---a BART-large generator (400M parameters) augmented with a Dense Passage Retriever (DPR) over a FAISS index of 21 million Wikipedia passages---in two variants: RAG-Sequence (one retrieved passage conditions the whole output) and RAG-Token (a different passage per token). RAG set state-of-the-art on three open-domain QA benchmarks (44.5 EM on Natural Questions, 45.5 on WebQuestions, 52.2 on CuratedTrec) and was judged more factual than BART in 42.7\,\% of human comparisons versus 7.1\,\%, with more diverse outputs (distinct tri-gram ratio 53.8\,\% vs.\ 32.4\,\%); critically, its retrieval index can be hot-swapped without retraining. The retrieval component builds on \cite{karpukhin_dense_2020}, whose dual-encoder DPR outperforms BM25 by 9--19\,\% in top-20 retrieval accuracy, and \cite{reimers_sentence-bert_2019}'s Sentence-BERT, which adapts BERT via a Siamese training objective for efficient cosine-similarity comparison---the \texttt{sentence-transformers} library built on this work provides Training Copilot's embedding model (\texttt{all-MiniLM-L6-v2}, $\sim$90\,MB, running locally without an embedding API). \cite{gao_retrieval-augmented_2024} survey the progression from Na\"ive RAG (retrieve-then-generate) through Advanced RAG (query rewriting, re-ranking) to Modular RAG (pluggable components); Training Copilot implements the na\"ive form---embed, retrieve by cosine similarity, generate with retrieved context---sufficient for its curated corpus of $\sim$3{,}000 chunks from eight public-domain books.

\subsection{AI-Assisted Sports Coaching}\label{sec:sota:coaching}

AI-assisted sports coaching is an emerging field where three limitations recur: insufficient access to the athlete's actual data, limited cross-domain reasoning, and single-agent architectures that do not scale to heterogeneous tool sets.

On the data-access and personalisation gap, \cite{monteiro-guerra_personalization_2020} reviewed 28 papers covering 17 real-time physical-activity coaching apps across seven personalisation dimensions: feedback was universal (17/17) and goal setting near-universal (15/17), but context-awareness appeared in only 3/17 and behavioural adaptation in only 2/17, with implementation details given for just one context-aware system. \cite{luo_promoting_2021} reach a similar conclusion for conversational agents specifically---chatbots sustain engagement, but most lack wearable-data access---and \cite{puce_harnessing_2025} confirm this in a systematic review of 10 studies on generative AI for exercise prescription: AI-generated plans adhere to foundational training principles but lack specificity, progression, and, most critically, real-time physiological feedback, with identical prompts sometimes yielding different outputs.

On architecture, the most directly relevant prior system is GPTCoach~\citep{joerke_gptcoach_2024}, an LLM-based coaching chatbot evaluated at CHI~2025. It pairs GPT-4 with a prompt-chaining design (separate calls for dialogue-state classification, motivational-interviewing strategy, and tool-call decisions) and two tools querying three months of Apple HealthKit data via HL7 FHIR; in a 16-participant lab study it scored 4.8/5 on feeling supported and 4.5/5 on empathy, with 93\,\% MI-consistent behaviour by external coding. Yet GPTCoach is a single agent with two tools and one data source, integrates no weather, calendar, or route data, and was evaluated only on a single onboarding session. At the architectural level, \cite{vemuri_multiagent_2021} proposed SMART-ALEC, a multi-agent health-coaching system with six modules per user agent (planner, scheduler, learner, context analyser, executor, communicator), prototyped as BeSmart (SMS-based, beta-tested with 250 agents) and a JIT intervention system on Apple Watch with nightly-updated decision trees; predating MCP/A2A and coordinating via ad-hoc GPGP-style commitment exchange, it nonetheless validates the modular-specialist-with-central-coordinator philosophy that Training Copilot extends with standardised protocols and LLM-based reasoning in place of finite-state dialogues.

On the human--AI boundary, \cite{barger_artificial_2025} ran a mixed-methods RCT ($N=52$) comparing sessions where participants believed they were coached by AI (actually a human behind a live-motion avatar) versus openly by a human: working-alliance scores did not differ significantly ($M=72.73$ vs.\ $74.50$; $t(50)=-0.71$, $p=0.48$), and participants in the AI condition used fewer impression-management behaviours and disclosed more sensitive topics---supporting a synergy model of narrowly focused AI coaches working alongside humans. \cite{gabarron_human_2024} survey AI-driven solutions for physical activity more broadly, finding recommender systems most studied, followed by conversational agents, with sparse long-term evidence for either.

In summary, prior coaching systems have either data access (wearable-native apps, limited to one source) or reasoning capability (general-purpose LLMs, without live data), but none combine multi-source integration, cross-domain reasoning, and literature-grounded advice via standardised protocols.

\subsection{Summary and Research Gap}\label{sec:sota:summary}

Each research stream reviewed above has matured individually, yet all share one limitation: isolation from the complementary capabilities integrated training support requires. Agentic AI has shown that domain-scoped multi-agent architectures outperform monolithic designs, but relies on ad-hoc registries rather than standardised protocols; MCP and A2A solve the resulting $M \times N$ integration problem in principle but have not been combined in a domain-specific application; sports science has established that training decisions require synthesising load, recovery, and context, yet consumer platforms display these signals without reasoning across them; AI-coaching systems demonstrate that LLM-based guidance is perceived as personalised and actionable but trade data access off against reasoning capability; and RAG enables grounding in verifiable sources but has not been applied to sports coaching alongside live physiological data.

Table~\ref{tab:sota_comparison} makes this concrete, scoring the most directly comparable prior systems---academic multi-agent frameworks, the closest published MCP extension, and the leading commercial AI coaching products---against four capabilities (multi-agent coordination, standardised protocols, multi-source live data, retrieval-grounded knowledge) and the application domain each targets. No prior system combines all four, and none of the few sports/health systems combine more than one.

\begin{table}[ht]
\begin{center}
\begin{minipage}{\textwidth}
\caption[Comparison of prior agentic and coaching systems]{Comparison of prior agentic and coaching systems against the research-gap dimensions identified in this section. \checkmark\ = demonstrated; $\times$ = not demonstrated or not present.}\label{tab:sota_comparison}
\begin{tabularx}{\textwidth}{p{2.9cm}*{5}{>{\centering\arraybackslash}X}}
\toprule
\textbf{System} & \textbf{Multi-Agent} & \textbf{Standardised Protocol} & \textbf{Multi-Source Live Data} & \textbf{RAG-Grounded} & \textbf{Sports/Health Domain} \\
\midrule
HuggingGPT~\citep{shen_hugginggpt_2023} & \checkmark & $\times$ (ad-hoc) & $\times$ (model hub) & $\times$ & $\times$ \\
\addlinespace
AutoGen~\citep{wu_autogen_2023} & \checkmark & $\times$ (ad-hoc) & $\times$ (generic) & $\times$ & $\times$ \\
\addlinespace
CA-MCP~\citep{jayanti_enhancing_2026} & $\times$ (1 LLM + reactive servers) & \checkmark (MCP extension) & $\times$ (benchmark sandbox) & $\times$ & $\times$ (travel/logistics) \\
\addlinespace
GPTCoach~\citep{joerke_gptcoach_2024} & $\times$ (single agent) & $\times$ & $\times$ (HealthKit only) & $\times$ & \checkmark \\
\addlinespace
SMART-ALEC~\citep{vemuri_multiagent_2021} & \checkmark (per-user, ad-hoc GPGP) & $\times$ & $\times$ (SMS/watch only) & $\times$ & \checkmark \\
\addlinespace
WHOOP Coach~\citep{whoop_coach_2024} & $\times$ & $\times$ & $\times$ (WHOOP only) & $\times$ & \checkmark \\
\addlinespace
Athletica.ai~\citep{athletica_ai_2024} & $\times$ & $\times$ & \checkmark (Garmin + Strava) & $\times$ & \checkmark \\
\addlinespace
General-purpose LLM~\citep{puce_harnessing_2025} & $\times$ & $\times$ & $\times$ (no API access) & $\times$ & $\times$ (generic) \\
\addlinespace
\textbf{Training Copilot (this work)} & \checkmark & \checkmark (MCP + A2A) & \checkmark (5 live sources) & \checkmark & \checkmark \\
\botrule
\end{tabularx}
\end{minipage}
\end{center}
\end{table}

The research gap is therefore not in any individual capability but in their integration: no existing system combines standardised multi-source access (MCP), multi-agent coordination (A2A), domain-scoped specialist reasoning (LangGraph ReAct), and retrieval-grounded knowledge (RAG) for athlete-facing training support. Training Copilot's architecture maps each limitation to a specific design decision, detailed in Section~\ref{sec:architecture}.
